\section{Contratos OpenAPI y trazabilidad de interfaz}
\label{sec:contratos-openapi}

La E5 incorpora contratos OpenAPI 3.1.0 versionados para describir de forma
reproducible la superficie HTTP de los cuatro servicios de dominio y de la
fachada expuesta por el API Gateway. Los contratos se conservan en
\path{docs/openapi/}:

\begin{itemize}
  \item \texttt{auth-service-openapi.json};
  \item \texttt{usuarios-service-openapi.json};
  \item \texttt{academico-laboratorios-service-openapi.json};
  \item \texttt{reservas-solicitudes-service-openapi.json}; y
  \item \texttt{api-gateway-openapi.json}.
\end{itemize}

\subsection{Cobertura y fachada}

El validador contrasta dinámicamente las anotaciones de los controladores con
cada snapshot. Auth cubre 12/12 operaciones, Usuarios 44/44, Académico 71/71
y Reservas 79/79, sin operaciones adicionales. La fachada del Gateway contiene
159 rutas y 224/224 operaciones; también verifica las 6/6 familias configuradas
en \texttt{GatewayRoutes.java}. Este contrato representa rutas canónicas y los
alias heredados que el Gateway todavía expone, sin atribuirle controladores
propios.

\subsection{Seguridad contractual}

Las operaciones externas protegidas declaran un esquema HTTP Bearer con formato
JWT. Los endpoints públicos de autenticación no exigen ese esquema. Las APIs
internas de los servicios usan un esquema \texttt{apiKey} en el encabezado
\texttt{X-Internal-Api-Key}; esas rutas internas no forman parte del contrato de
fachada del Gateway.

\subsection{Producción, validación y CI}

Los snapshots se extraen determinísticamente de controladores, DTO y rutas
vigentes sin requerir que los servicios arranquen con infraestructura externa.
Desde la raíz se validan o regeneran con:

\begin{lstlisting}
python scripts/validar-contratos-openapi.py
python scripts/validar-contratos-openapi.py --generate
\end{lstlisting}

La validación comprueba JSON, estructura OpenAPI, identificadores de operación,
referencias locales, esquemas de seguridad y deriva método+ruta. El job de CI
\emph{Validate - OpenAPI contracts} ejecuta la validación sin dependencias
Python externas.

OpenAPI y Pact son complementarios. OpenAPI describe la superficie de la API,
sus parámetros y schemas; Pact verifica interacciones consumidor--proveedor
seleccionadas. Ninguno sustituye al otro. Como limitación, la extracción OpenAPI
es estática: mantiene trazabilidad estructural, pero no demuestra por sí sola
todos los comportamientos dinámicos ni las reglas de negocio durante ejecución.
